\section{引言：课程回顾与本讲主旨}

本讲是 MIT 6.S191 Introduction to Deep Learning 系列六节基础讲座的最后一讲。在之前的五讲中，课程系统地介绍了深度学习的核心方法论——从基础的全连接神经网络、序列建模（RNN / Transformer）、卷积神经网络（CNN），到生成模型（VAE / GAN）和深度强化学习。本讲将在回顾这些基础之上，聚焦两个核心议题：

\begin{enumerate}
    \item \textbf{深度学习的局限性}（Limitations）——泛化困境、对抗攻击、不确定性表征等；
    \item \textbf{新前沿}（New Frontiers）——Diffusion Models 和 Large Language Models (LLMs) 为代表的最新进展。
\end{enumerate}

\begin{figure}[H]
\centering
\includegraphics[width=0.85\textwidth]{slides-images/slide-01.jpg}
\caption{本讲主题：Deep Learning Limitations and New Frontiers\protect\footnotemark}
\end{figure}
\footnotetext{Slides 第1页。}

正如 Amini 在课上所说，深度学习已经在自动驾驶、医学影像、强化学习、生成建模、自然语言处理、金融、安全等众多领域带来了革命性进展。但作为技术从业者，我们不仅要理解这些算法的威力，更要清醒地认识其局限性，这样才能负责任地将 AI 部署到真实世界中。

\begin{figure}[H]
\centering
\includegraphics[width=0.85\textwidth]{slides-images/slide-10.jpg}
\caption{深度学习在各个领域的兴起\protect\footnotemark}
\end{figure}
\footnotetext{Slides 第10页。}

\subsection{本章小结}

本讲作为基础讲座的收官之作，承上启下：回顾深度学习的核心作为函数逼近器的本质，指出其局限性，并展望以 Diffusion Models 和 LLM 为代表的新前沿方向。

%% =====================================================
